\documentclass[11pt, a4paper]{article}

\usepackage{graphicx} % Required for inserting images
% \usepackage[utf8]{inputenc} % Quins caràcters es veuran
\usepackage[document]{ragged2e} % Posar els paràgrafs
\usepackage{amsmath,amsthm,amssymb,latexsym}
\usepackage{enumitem}
% \usepackage{xcolor}
\usepackage[x11names]{xcolor}
\usepackage{tikz}
\usepackage{tcolorbox}
\usepackage{tkz-euclide}
\usepackage{pgfplots}
% \usepackage{cancel}
\usepackage{hyperref}
\usepackage{mathrsfs,mathtools}
\usepackage{microtype}
\usepackage{fancyhdr}
\usepackage{parskip}
\usepackage{tkz-euclide}

\usepackage[catalan]{babel}

\renewcommand\refname{}
\usepackage[margin=1in]{geometry}



\pgfplotsset{compat=1.18}
\usetikzlibrary{patterns}
\usetikzlibrary{intersections}


\setlength{\headheight}{15.21004pt}

\DeclareMathOperator{\Ima}{Im}

%% Macros
\providecommand{\ol}{\overline}
\providecommand{\ul}{\underline}
\providecommand{\wt}{\widetilde}
\providecommand{\wh}{\widehat}
\providecommand{\eps}{\varepsilon}
\providecommand{\half}{\frac{1}{2}}
\providecommand{\inv}{^{-1}}
\newcommand{\dang}{\measuredangle} %% Directed angle
\providecommand{\CC}{\mathbb C}
\providecommand{\FF}{\mathbb F}
\providecommand{\NN}{\mathbb Z_{\geq 1}}
\providecommand{\QQ}{\mathbb Q}
\providecommand{\RR}{\mathbb R}
\providecommand{\PP}{\mathbb P}
\providecommand{\ZZ}{\mathbb Z}
\renewcommand{\AA}{\mathbb A}
\providecommand{\dd}{\mathrm d}
\newcommand{\mfrk}{\mathfrak}
\newcommand{\KK}{\mathbb{K}}
\newcommand{\TT}{\mathcal{T}}
\providecommand{\ts}{\textsuperscript}
\providecommand{\dg}{^\circ}
\providecommand{\ii}{\item}
\DeclareMathOperator*{\lcm}{lcm}
\DeclareMathOperator*{\argmin}{arg min}
\DeclareMathOperator*{\argmax}{arg max}
\DeclareMathOperator*{\rad}{rad}
\DeclareMathOperator*{\spec}{Spec}
\DeclareMathOperator*{\mult}{mult}
\providecommand{\bigO}{\mathcal O}




% theorem environments
% Les versions estrellades no estan numerades 
\theoremstyle{definition}
    \newtheorem{theorem}{Teorema}
    \newtheorem{lemma}[theorem]{Lema}
    \newtheorem{claim}[theorem]{Claim}
    \newtheorem{corollary}[theorem]{Corol·lari}
    \newtheorem{definition}[theorem]{Definició}
    \newtheorem{proposition}[theorem]{Proposició}
    \newtheorem{example}[theorem]{Exemple}
    
    \newtheorem*{theorem*}{Teorema}
    \newtheorem*{lemma*}{Lema}
    \newtheorem*{claim*}{Claim}
    \newtheorem*{corollary*}{corol·lari}
    \newtheorem*{definition*}{Definició}
    \newtheorem*{proposition*}{Proposició}
    \newtheorem*{example*}{Exemple}
    \newtheorem*{proof*}{Demostració}
    
\theoremstyle{remark}
    \newtheorem{remark}[theorem]{Remarca}
    \newtheorem*{remark*}{Remarca}

\usepackage{hyperref}
\renewcommand{\thesubsection}{}
\renewcommand{\thesection}{}

\title{Introducció a l'àlgebra commutativa}
\author{Bernat Esteve Sagarra}

\begin{document}
\pagestyle{fancy}
% 
\lhead{Bernat Esteve}
\rhead{Laboratori 3}
\begin{tcolorbox}[title=\section{Exercici 1}]
Segons el teorema demostrat a teoria, si $F$ és una corba irreductible de grau $d$, es compleix que la suma de multiplicitats sobre els seus punts satisfà la següent fita:
\[\sum_{p \in V_+(F)} m_p(F)(m_p(F)-1) \le (d-1)(d-2)\]
\begin{enumerate}
    \item Utilitzeu aquest resultat per demostrar que una quàrtica ($d=4$) irreductible pot tenir, com a màxim, 3 singularitats de multiplicitat 2.
    \item Demostreu que no existeix cap corba irreductible de grau 4 amb quatre punts $p_i$ de multiplicitat $m_i=2$ per a $i=1,...,4$.
    \item Malgrat l'apartat anterior, si la corba no és irreductible, sí que pot existir. Doneu un exemple explícit d'una quàrtica (que no sigui irreductible però sí reduïda) que tingui exactament 4 singularitats de multiplicitat 2. Trobeu-ne els punts i descriviu-ne els cons tangents.
\end{enumerate}
\end{tcolorbox}
\subsection{Exercici 1}
Notem que si tinguéssim 4 singularitats de multiplicitat 2, aleshores
\[
    8=2+2+2+2>3\times 2=6
\]
Que contradiu el teorema.\par\vspace{5mm}
Però notem que si considerem la següent corba singular:
\[
    F=xyz(x+y+z)
\]
Això té ${4\choose 2}=6$ singularitats.\par
Les singularitats són en el 
\[
x^2yz+xy^2z+xyz^2
\]
Per tant, $[0:0:1]$, $[0:1:0]$, $[1:0:0]$, $[1:-1:0]$, $[0:1:-1]$ i $[-1:0:1]$. Que tenen con tangent les rectes $x=0$, $y=0$, $z=0$, $x+y+z=0$.
\newpage
\begin{tcolorbox}[title=\section{Exercici 2}]
Siguin $F$ i $G$ dues corbes projectives irreductibles de graus $m$ i $n$ respectivament, sense components en comú. Sabem que per al producte de corbes es compleix que $\text{Sing}(F \cdot G) = \text{Sing}(F) \cup \text{Sing}(G) \cup (V_+(F) \cap V_+(G))$.
\begin{enumerate}
    \item Deduiu una fita superior per al nombre total de singularitats de la corba $C = F \cdot G$ exclusivament en funció dels graus $m$ i $n$. Indiqueu de quina manera intervé el Teorema de Bézout en la vostra fita.
    \item Apliqueu la fita obtinguda al cas en què $F$ i $G$ són dues còniques irreductibles i llises (no singulars) diferents. Quin és el nombre màxim de singularitats que pot tenir $F \cdot G$?
    \item Suposeu ara que $m=3$ (cúbica llisa) i $n=1$ (recta). Sota quina condició geomètrica la corba resultant tindrà exactament 2 singularitats?
\end{enumerate}
\end{tcolorbox}
\subsection{Exercici 2}
\[
\sum_{p \in V_+(F\cdot G)} m_p(F)(m_p(F)-1)\le (m-1)(m-2)+(n-1)(n-2)-mn
\]
Utilitzant aquesta fita, tenim que el màxim nombre de singularitats és 4 (només les de $V_+(F)\cup V_+(G)$).\par
Si la cúbica és llisa, aleshores els punts singulars són les interseccions: i per tant, hi ha 2 singularitats si la recta és una recta tangent qualsevol (de multiplicitat 2).
\newpage
\begin{tcolorbox}[title=\section{Exercici 3}]
Sigui $F$ una corba de Plücker de grau $d$ amb $n$ nodes i $c$ cúspides. A teoria s'ha definit $d^*$ com el nombre de tangents a $F$ en punts no singulars que passen per un punt $q$ exterior general (que no pertany a cap recta tangent en un punt singular), de manera que $d^* = d(d-1) - 2n - 3c$ segons la primera fórmula de Plücker.
\begin{enumerate}
    \item Calculeu el valor de $d^*$ per a una cúbica no singular.
    \item Sabent que el grau de la corba dual $F^*$ coincideix amb $d^*$, quin serà el grau de la corba dual d'aquesta cúbica no singular?
    \item Comproveu com canvia la situació per a una quàrtica: calculeu quantes tangents passen per un punt general si la quàrtica és completament llisa ($n=0, c=0$) i compareu-ho amb una quàrtica irreductible que assoleixi el màxim nombre de nodes permès teòricament (segons el Problema 1).
\end{enumerate}
\end{tcolorbox}
\subsection{Exercici 3}
\newpage
\begin{tcolorbox}[title=\section{Exercici 4}]
Considerem la família de corbes afins donades per $f_k(x,y)=y^2-x^k$ amb $k\geq2$.
\begin{enumerate}
    \item Trobeu els punts singulars a l'espai afí per a qualsevol valor de $k$.
    
    \item Homogeneïtzeu el polinomi per obtenir la corba projectiva $F_k(x,y,z)$ i estudieu si hi ha punts singulars addicionals a la recta de l'infinit.
    
    \item Centrant-nos en l'origen afí, classifiqueu la singularitat geomètrica per als casos $k=2$, $k=3$ i $k=4$. Determineu explícitament els cons tangents per identificar si es tracta d'un node (singularitat ordinària de multiplicitat 2), una cúspide, o un altre tipus.
\end{enumerate}

\end{tcolorbox}
\subsection{Exercici XXX}

\newpage
\begin{tcolorbox}[title=\section{Exercici XXX}]
Considereu la corba projectiva definida pel polinomi de grau 4:

\[
F(x,y,z)=(x^2 +y^2-z^2)(x ^2+y ^2-2zx)
\]
\begin{enumerate}
    \item Trobeu els factors irreductibles de $F$ i demostreu que de manera independent defineixen corbes llises (no singulars).
    
    \item Trobeu el conjunt complert de punts singulars $\mathrm{Sing}(F)$. \textit{Indicació: Podeu utilitzar el mètode de descomposició descrit en el Problema 2 si us resulta més àgil que derivar l'expressió completa.}
    
    \item Per a cadascun dels punts singulars de l'apartat anterior, trobeu-ne les rectes tangents i determineu si la singularitat és ordinària o no. Com es tradueix aquest fet geomètricament respecte a l'angle d'intersecció de les components irreductibles?
\end{enumerate}
\end{tcolorbox}
\subsection{Exercici 5}
\end{document}
% ---
% \newpage
% \begin{tcolorbox}[title=\section{Exercici XXX}]
% \end{tcolorbox}
% \subsection{Exercici XXX}
